coupled in the eigenequations (14) only by means of the squared eigenfrequency $\omega^2$.
This permits the existence of purely real, standing-wave eigenfunctions.

The eigenequations (14) comprise a fifth-order differential system; for any given
choice of the complex frequency, $\omega$, there are five linearly independent, complex solutions
to equations (14). Physical boundary conditions outlined in the next two sections
make four of the five solutions unacceptable. The remaining solution is physically acceptable for all choices of $\omega$, providing one permits arbitrary admixtures of ingoing and
outgoing radiation. The restriction to purely outgoing radiation restricts $\omega$ to take on a
discrete set of values, as we shall see in \S~IIIc.

\subsection*{b) Boundary Conditions at the Star's Center and Surface}

By expanding the eigenequations (14) about $r = 0$ one finds that only two of the
five independent solutions are physically acceptable at the center of the star. (See
Appendix~D for analysis.) The two acceptable solutions have the form:
\begin{equation}
  K = Ar^l + \cdots, \qquad
  H = Ar^l + \cdots, \qquad
  W = Br^{l+1} + \cdots,
\label{eq:15a}
\end{equation}
\begin{equation}
  V = -(B/l)r^l + \cdots, \qquad \text{near } r = 0\;.
\notag
\end{equation}

Here $A$ and $B$ are independent, freely specifiable constants. These acceptable solutions
are characterized by small fluid motions and by small perturbations in the density, the
pressure, and the geometry of spacetime. The unacceptable solutions---which are discussed in Appendix~D---are characterized by perturbations in the density, the pressure,
and/or the geometry of spacetime which diverge as $r = 0$.

An expansion of the eigenequations (14) about $r = R$ (surface of star) reveals that
four of the five independent solutions are physically acceptable there. (See Appendix~D
for analysis.) The acceptable solutions have the following expansion upon $(R - r)$:
\begin{align}
  K &= k_0 + k_1(R-r) + \cdots, &
  H &= h_0 + h_1(R-r) + \cdots, \notag \\
  W &= w_0 + w_1(R-r) + \cdots, &
  V &= v_0 + v_1(R-r) + \cdots,
\label{eq:15b}
\end{align}
near $r = R$.

Of the constants which enter into this expansion only four are freely specifiable. The
remaining constants are fixed by the asymptotic form of the eigenequations and by the
demand that the Lagrangian perturbation in the pressure vanish at the star's surface:
\begin{equation}
  \Delta p \equiv \lim_{r \to R^-} \gamma p \left[
    \frac{e^{-\lambda/2}}{r^2} W' - \frac{l(l+1)}{r} V
    + \frac{1}{2} H_1 + K
  \right] P_l(\cos\theta) = 0\;.
\label{eq:16}
\end{equation}

For example, if the pressure-density relation near the star's surface is polytropic and the
adiabatic index behaves smoothly,
\begin{equation}
  p = \rho^{1+1/n} + \cdots, \qquad \mu = \bar{\mu}(R-r)^n + \cdots,
\label{eq:17a}
\end{equation}
\begin{equation}
  \gamma = \gamma_1 + \gamma_2(R-r) + \cdots,
\notag
\end{equation}
then the four constants $k_0$, $h_0$, $w_0$ and $w_0$ are freely specifiable in the acceptable solutions
(15b). The eigenequation (14d) fixes $v_0$ in terms of these four arbitrary constants
\begin{equation}
  \omega^2 R(R - 2M)^{-1} v_0 = R^{-1}(1 - 2M/R)^{1/2} (N+1) w_0\;,
\label{eq:17b}
\end{equation}
and the eigenequations (14a)--(14d) all conspire to fix the remaining constants in the
ascending power series. Here $M \equiv m(r = R)$ is the total mass of the star. (In this paper
$M$ is used both as the star's total mass and as the spherical harmonic projection index.)

In general, the one physically unacceptable solution at the star's surface is that
